\SourceLeaf{009}
\Article{2}
What has been set out so far is adequately proved in books on algebra and nowhere offends geometric rigour. But analysts seem to have adopted too hastily, and without a prior sound proof, the theorem on which almost the entire theory of equations rests: \textit{every function such as $X$ can always be resolved into $m$ linear factors}, or, what agrees with this completely, \textit{every equation of degree $m$ actually has $m$ roots}. Since even quadratic equations very often lead to cases that contradict this theorem, algebraists, in order to bring these cases under it, were forced to invent an imaginary quantity whose square is $-1$. They then recognized that, if quantities of the form $a+b\sqrt{-1}$ are admitted in the same way as real quantities, the theorem holds not only for quadratic equations but also for cubic and biquadratic equations. Yet this by no means justified concluding that, once quantities of the form $a+b\sqrt{-1}$ are admitted, every equation of the fifth or a higher degree can be satisfied, or, as it is usually expressed (though I would not favour this slippery wording), that the roots of every equation can be reduced to the form $a+b\sqrt{-1}$. In substance, this theorem differs in no respect from the one stated in the title of this work; giving a new rigorous proof of it is the aim of the present dissertation.

Moreover, ever since analysts discovered that there are infinitely many equations that would have no root at all unless quantities of the form $a+b\sqrt{-1}$ were admitted, such fictitious quantities have been regarded as a special kind of quantity, called \textit{imaginary} to distinguish them from \textit{real} quantities, and have been introduced throughout analysis. By what right? I do not discuss that here. --- I shall complete my proof without any help from imaginary quantities, although I too would be entitled to use the same freedom that all recent analysts have used.

\Article{3}
Although what most elementary books offer as a proof of our theorem is so flimsy and so far removed from geometric rigour as scarcely to deserve mention, I shall nevertheless touch on it briefly so that nothing is left out. `To show that every equation
\[
x^m+Ax^{m-1}+Bx^{m-2}+\text{etc.}+M=0
\]
\SourceLeaf{010}
or $X=0$ actually has $m$ roots, they undertake to prove that $X$ can be resolved into $m$ linear factors. For this purpose they assume $m$ linear factors $x-\alpha$, $x-\beta$, $x-\gamma$, etc., where $\alpha$, $\beta$, $\gamma$, etc. are still unknown, and set their product equal to the function $X$. Then, by comparing coefficients, they obtain $m$ equations from which, they say, the unknowns $\alpha$, $\beta$, $\gamma$, etc. can be determined, since their number is also $m$. Namely, $m-1$ unknowns can be eliminated, leaving an equation containing only whichever unknown one wishes.' Passing over the other objections that could be made to such reasoning, I ask only: how can we be certain that this final equation actually has any root? Why could it not happen that no magnitude in the entire range of real and imaginary quantities satisfies either this final equation or the proposed one? --- Experts will, moreover, readily see that this final equation must be \textit{entirely identical} with the proposed one if the calculation has been carried out correctly: namely, eliminating the unknowns $\beta$, $\gamma$, etc. must give the equation
\[
\alpha^m+A\alpha^{m-1}+B\alpha^{m-2}+\text{etc.}+M=0
\]
There is no need to say more about this reasoning.

Some authors who seem to have noticed the weakness of this method assume as an \textit{axiom} that every equation actually has roots: if not possible roots, then impossible ones. They do not seem to have made sufficiently clear what they mean by possible and impossible quantities. If possible quantities are to mean the same as real quantities, and impossible quantities the same as imaginary ones, that axiom can by no means be admitted: it necessarily requires a proof. Yet the expressions do not seem intended in that sense; rather, the axiom seems to mean this: `Although we are not yet certain that there must be $m$ real or imaginary quantities satisfying a given equation of degree $m$, we shall nevertheless suppose this for a while; for if it should happen that so many real and imaginary quantities cannot be found, we shall certainly have a way out by saying that the remaining ones are impossible.' If someone prefers this wording to saying simply that the equation would not have so many roots in this case, I have no objection. But if he then uses these impossible roots as though they were something real, and says, for example, that the sum of all the roots of the equation $x^m+Ax^{m-1}+\text{etc.}=0$ is $=-A$, even if there are impossible roots among them (an expression that properly means \textit{even if some are missing}), I cannot approve of this at all. For impossible roots\SourceLeaf{011}, understood in this sense, are still roots; that axiom therefore cannot possibly be admitted without proof, and would it be unreasonable to wonder whether there might be equations that do not even have impossible roots?\HistoricalNote{1}

\Article{4}
Before reviewing other mathematicians’ proofs of our theorem and setting out what seems objectionable to me in each, I note that it is enough to show only that every equation of any degree
\[
x^m+Ax^{m-1}+Bx^{m-2}+\text{etc.}+M=0
\]
or $X=0$ (where the coefficients $A,B$, etc. are supposed real) can be satisfied in at least one way by a value of $x$ of the form $a+b\sqrt{-1}$. For it is known that $X$ will then be divisible by the real quadratic factor $xx-2ax+aa+bb$ if $b$ is not $=0$, and by the real linear factor $x-a$ if $b=0$. In either case the quotient will be real and of lower degree than $X$; and since, by the same reasoning, this quotient must have a real factor of the first or second degree\SourceLeaf{012}, it is evident that continuing this operation will ultimately resolve the function $X$ into real linear or quadratic factors, or, if one prefers to use two imaginary linear factors in place of each real quadratic factor, into $m$ linear factors.
